\section{Results}
\label{sec:results}

Every number below resolves to an entry in the evidence ledger, and every table and figure is
generated from it rather than transcribed.

\subsection{The channel is real, and it is capped by the construction}

Table~\ref{tab:capacity} and Figure~\ref{fig:capacity} give the realized watermark information per DNA
base for each policy: $\capacityLow{}$ to $\capacityHigh{}$ bit per base against a structural ceiling
of $\capacityCeiling{}$. The ceiling is not a limit of these models. Partition maximal coupling
couples exactly one bit into each 6-mer, so $1/6$ bit per base is the most the construction can carry
no matter how much entropy the model supplies, and the measured values sit just under it. The nominal
alphabet of $4{,}096$ tokens is therefore not the quantity that matters, and reporting it would
overstate the channel by more than an order of magnitude.

% GENERATED FILE -- do not edit. Regenerate with paper/scripts/make_tables.py
\begin{table}[t]
  \centering
  \small
  \begin{tabular}{lrrrr}
    \toprule
    Policy & bit/base & 95\% interval & states & prompts \\
    \midrule
    \policyCtok & 0.1540 & [0.1515, 0.1558] & 3072 & 24 \\
    \policyGtok & 0.1514 & [0.1398, 0.1578] & 3072 & 24 \\
    \policyGbp & 0.1574 & [0.1548, 0.1593] & 3072 & 24 \\
    \bottomrule
  \end{tabular}
  \caption{Realized maximal-coupling watermark information per DNA base for each released generation policy, with equal-weight prompt-cluster percentile intervals. The construction's structural ceiling is one coupled bit per 6-mer token, or $1/6$ bit per base. Sequential states within a prompt are not independent samples, which is why the interval resamples prompts rather than states.}
  \label{tab:capacity}
\end{table}


\subsection{Clean detection is complete from the floor of the search}

Table~\ref{tab:clean-detection} and Figure~\ref{fig:clean-detection} report clean detection at a
calibrated false-positive rate of $\targetFPR{}$. Every watermarked sequence is detected at every
evaluated length down to $\shortestDetectedBases{}$ bases, which is the shortest sequence the declared
search will score rather than a measured limit: the search declares a minimum window of $16$ tokens.

A flat detection rate hides the quantity that actually varies, so we admit the separation margin
separately. At $\shortestDetectedBases{}$ bases the weakest watermarked statistic and the strongest
pooled null statistic touch; from $\strictSeparationBases{}$ bases onward the watermarked statistics
sit strictly above every null. Figure~\ref{fig:clean-detection} shows the watermarked statistic growing
with the square root of the length while the null maximum stays flat, which is the whole mechanism in
one panel.

% GENERATED FILE -- do not edit. Regenerate with paper/scripts/make_tables.py
\begin{table}[t]
  \centering
  \small
  \begin{tabular}{lrrrrr}
    \toprule
    Policy & fully detected (b) & strict separation (b) & threshold & achieved FPR & TPR \\
    \midrule
    \policyCtok & 96 & 144 & 4.066 & 0.00937 & 1.000 \\
    \policyGtok & 96 & 144 & 3.760 & 0.00625 & 1.000 \\
    \policyGbp & 96 & 96 & 3.536 & 0.00625 & 1.000 \\
    \bottomrule
  \end{tabular}
  \caption{Clean detection. The shortest fully detected length is the shortest sequence the declared search will score at all, so it is a floor of the experiment rather than a measured limit; the strict-separation length is where the weakest watermarked statistic passes the strongest pooled null. The threshold, achieved false-positive rate, and detection rate are quoted at 3{,}072 bases and come from null trials that repeat the identical declared search.}
  \label{tab:clean-detection}
\end{table}


\subsection{Substitutions are survivable; synchronization is the binding limit}

Table~\ref{tab:edit-channel} collects the edit channel. Substitution tolerance is high: every sequence
is still detected at $\substitutionTolerance{}$ substitutions per base on the least tolerant policy
and at $\substitutionToleranceHigh{}$ on the other two (Figure~\ref{fig:substitution}). At roughly one
base in five replaced, the mark still comes through on every sequence.

Insertions and deletions are different in kind rather than degree. Deleting one base moves every later
base one position left, so every 6-mer after the deletion is cut at the wrong boundary and the
remainder of the sequence leaves the reading grid. The unwindowed detector tolerates only
$\deletionTolerance{}$ deletions per base, one to two orders of magnitude below its substitution
tolerance. A declared sliding-window search recovers a factor of several, to
$\windowedDeletionTolerance{}$, and then meets a second wall that is information-theoretic rather than
a search failure: below some number of consecutive intact bases there is no window long enough to
carry a detectable signal, however cleverly it is placed (Figure~\ref{fig:indel}).

Two predictions were wrong here in ways worth recording. The first underestimated detection badly
because the construction resynchronizes for free: after $k$ deleted bases, phase $k \bmod 6$ combined
with key offset $\lceil k/6 \rceil$ recovers the original grid, so the declared phase and offset search
already absorbs tens of bases of drift. The second used an unsigned drift range, which cannot reach any
post-insertion segment at all --- a silent total failure on the insertion channel rather than a loss of
power. Both corrections are recorded in the frozen protocols rather than smoothed away.

% GENERATED FILE -- do not edit. Regenerate with paper/scripts/make_tables.py
\begin{table}[t]
  \centering
  \small
  \begin{tabular}{lrrrr}
    \toprule
    Policy & substitution & deletion & insertion & deletion, windowed \\
    \midrule
    \policyCtok & 0.15 & 0.002 & 0.002 & 0.005 \\
    \policyGtok & 0.2 & 0.001 & 0.001 & 0.005 \\
    \policyGbp & 0.2 & 0.001 & 0.002 & 0.01 \\
    \bottomrule
  \end{tabular}
  \caption{The edit channel: largest per-base edit rate at which every sequence is still detected. Substitutions are quoted at 3{,}072 bases and the indel channels at 1{,}536 bases, the longest length each experiment evaluated. Insertions and deletions defeat the unwindowed detector one to two orders of magnitude below the substitution tolerance, because one inserted or deleted base moves every later token off the reading grid. A declared sliding-window search moves the deletion limit but meets a second, information-theoretic wall.}
  \label{tab:edit-channel}
\end{table}


Crops, unknown phase, and reverse complementation are handled by the declared search
(Figure~\ref{fig:crop-strand}). Every condition that fails is one whose required key-stream offset lies
outside the declared search, which is a property of the search rather than of the model, and the three
policies fail on exactly the same conditions. Widening the search recovers all but the longest crop and
raises the calibrated threshold; both effects are inside the calibration because the nulls repeat the
widened search.

\subsection{Measured against the published alternatives, our construction loses}

Table~\ref{tab:baselines} and Figure~\ref{fig:baselines} compare partition coupling with inverse
transform and exponential sampling on the same policies, prompts, lengths, and control, each
calibrated on its own nulls at the same target rate and searching the identical hypothesis set.

The preregistered headline was the shortest length at which each method detects every sequence. It is
$\baselineFullyDetectedBases{}$ bases for all three methods on all three policies --- the floor of the
search --- so that quantity cannot separate them, and we report it as the null result it is rather than
as agreement.

What separates them is the signal each construction puts into one token, measured in units of its own
null. Exponential sampling carries at least $\expSignalRatio{}$ times partition coupling's per-token
signal and inverse transform at least $\itsSignalRatio{}$ times. Partition coupling is not
underperforming: it is at its one-bit-per-token ceiling, and the measured value equals $2p-1$ for the
generator-side agreement rate $p$. Exponential sampling reads a fine-grained keyed value attached to
the emitted token and is not bounded by one bit, and that is the entire difference.

% GENERATED FILE -- do not edit. Regenerate with paper/scripts/make_tables.py
\begin{table}[t]
  \centering
  \small
  \begin{tabular}{lrrrrr}
    \toprule
    Policy & ours & inverse transform & exponential & fully detected (b) & natural-DNA exceedance \\
    \midrule
    \policyCtok & 0.94--0.97 & 1.38--1.39 & 7.19--7.86 & 96 & 0.069 \\
    \policyGtok & 0.98--1.00 & 1.38--1.39 & 7.43--8.02 & 96 & 0.044 \\
    \policyGbp & 0.99--1.00 & 1.39--1.39 & 7.29--7.62 & 96 & 0.019 \\
    \bottomrule
  \end{tabular}
  \caption{The three exact-marginal constructions on one corpus and one calibrated detector, each calibrated on its own nulls at the same target rate and searching the identical hypothesis set. Signal per token is standardized in units of each method's own null, which is the only sense in which the three are comparable. Detection rate is not tabulated because it is $1.000$ for all three methods at every evaluated length, so the shortest fully detected length is identical and uninformative. The last column is the largest rate at which natural public DNA under a wrong key exceeds a threshold calibrated on model-generated nulls.}
  \label{tab:baselines}
\end{table}


The comparison also produced a finding we were not looking for. Thresholds are calibrated on
model-generated nulls, but a deployed verifier is handed natural DNA. Natural public DNA under wrong
keys exceeds those thresholds at up to $\naturalDNAExceedance{}$, above the $\targetFPR{}$ target, and
the two published baselines are more exposed to this than partition coupling at short lengths.
Calibration does not transfer from generated to natural sequences, and it does not transfer equally
across constructions. This was visible only because public-DNA trials were reported separately and
never pooled into a threshold.

\subsection{Key reuse defeats the construction}

A deployed generator reuses its key, so an attacker accumulates outputs under one secret.
Table~\ref{tab:attacks} and Figure~\ref{fig:attacks} measure what that buys.

Take $\forgeryDonorsRequired{}$ watermarked outputs and, for each token position, copy that position's
6-mer from one of them. The result is a sequence the generator never produced, and the verifier accepts
it at rate $\forgeryDetectionRate{}$ at every length on every policy, with a statistic
indistinguishable from genuine output. Forgeries that happen to copy a single donor throughout are
rejected and redrawn, so no counted success is a verbatim copy.

This is structural, not a tuning failure. The detector's decision at read position $i$ depends only on
which half of the keyed partition at stream index $i$ the observed token falls into, and that partition
is a function of the key, the domain, and $i$ alone. The statistic is a sum of per-position,
content-blind indicators; nothing binds position $i$ to position $j$, and nothing binds the sequence to
a nonce, a length, or a payload. \textbf{The construction authenticates a distribution, not a
sequence.} Closing the gap requires a different construction, and the same argument applies to any
position-indexed, content-blind keyed statistic, which includes both baselines --- though we measured
it only on ours.

Removal is the mirror image. Permuting 6-mer positions moves each token to a position whose keyed split
is unrelated to it, so each check becomes a fair coin: a full positional shuffle drops detection to
$\shuffleDetectionRate{}$ everywhere. A shuffle bounded to blocks of two tokens retains about half the
statistic, and because the statistic grows with the square root of the length, half of a large statistic
still clears the threshold; detection after a pairwise shuffle therefore \emph{rises} with length, and
removal requires wider blocks as sequences grow.

% GENERATED FILE -- do not edit. Regenerate with paper/scripts/make_tables.py
\begin{table}[t]
  \centering
  \small
  \begin{tabular}{lrrrrrr}
    \toprule
    Policy & forgery accepted & donors & shuffled detected & removal width & proxy cost & reuse exceedance \\
    \midrule
    \policyCtok & 1.000 & 2 & 0.000 & 8 & 0.167 & 0.125 \\
    \policyGtok & 1.000 & 2 & 0.000 & 16 & 0.123 & 0.000 \\
    \policyGbp & 1.000 & 2 & 0.000 & 32 & 0.143 & 0.125 \\
    \bottomrule
  \end{tabular}
  \caption{Key reuse. \emph{Forgery accepted} is a vulnerability, not a result: a sequence spliced position-wise from watermarked outputs under the same key is accepted at every length, and two donor outputs suffice. \emph{Shuffled detected} is the detection rate after a full positional shuffle, so a low value is a successful removal; the next column is the narrowest block width whose shuffle removes the mark at 3{,}072 bases. The proxy cost is the largest relative change in any admitted sequence proxy, and it is small mainly because the attack preserves the 6-mer multiset exactly and the admitted proxies are composition statistics. The last column is one fixed key scored against sequences it did not generate.}
  \label{tab:attacks}
\end{table}


\subsection{Pricing the removal attack, and failing to}

A removal attack is only meaningful with its cost attached, so this deserved its own experiment. The
largest relative change in any admitted composition proxy under a full shuffle is
$\removalProxyCost{}$, on a single extreme-value statistic, with every other proxy moving far less.
That is a fact about the proxies rather than about the attack: the attack permutes whole 6-mers, so the
6-mer multiset is preserved \emph{exactly}, and a suite of counting statistics is blind to a
permutation of the units it counts.

We therefore implemented the two order-sensitive measures our own design specifies --- open reading
frame summaries and a score under an independent Markov reference fitted on held-out cohort prompts ---
and neither prices the attack.

The independent-model score is blind in practice: its largest relative shift across every removal and
spoofing condition is $\independentModelShift{}$. The reason is structural. A low-order Markov model
scores high-entropy generated DNA close to the uniform rate, and a 6-mer permutation leaves most short
contexts intact.

The reading-frame result is more instructive, because it first looked like a success. The relative
shift of the longest open reading frame runs from $0.07$ to $0.43$, larger than the composition shift
on most cells. An exact paired sign-flip test over the eight prompts dissolves it: of
$\orfConditionsTested{}$ removal conditions, at most $\orfDirectionalConditions{}$ shows a
directionally consistent change on any policy, and for every bounded block shuffle --- the attacker's
realistic option, and the widths that actually strip the mark --- the longest ORF rises about as often
as it falls. On high-entropy DNA a long reading frame is a chance extreme-value statistic, and
permuting the sequence re-rolls it rather than destroying a structure. A relative shift built from
absolute differences records the re-roll and hides the absence of direction
(Figure~\ref{fig:unpriced}).

So the limitation is sharpened rather than closed. We are not reporting that we neglected to price the
attack; we are reporting that we built both instruments our design named for the purpose and they
cannot see it. Pricing a bounded 6-mer rearrangement needs a measure of long-range order that none of
these three families supplies, and finding one is open work.

One asymmetry is worth noting, and it runs against the defender. The splice forgery disturbs the
composition proxies substantially more than any shuffle does, because splicing mixes DNA across
organisms. The attack that manufactures provenance is easier to spot by composition than the attack
that destroys it.

Finally, key reuse does not break calibration on the verifier's side: one fixed key scored against
sequences it did not generate exceeds its own threshold at a rate consistent with the target, within
the resolution that $8$ trials per cell allows.

\subsection{Feasibility}

Table~\ref{tab:runtime} gives the runtime envelope from single observations on one laptop under
uncontrolled desktop load: a feasibility statement, not a benchmark. Detection uses neither the model
nor the GPU. We report no peak-memory figure, because the recorded memory counters cannot support a
residency claim --- one driver counter exceeds the machine's physical memory --- and a number we cannot
defend is worse than none. All quoted times are mains-power measurements; the same generation runs about
ten times slower on battery, which is a reproducibility constraint rather than a property of the method.

% GENERATED FILE -- do not edit. Regenerate with paper/scripts/make_tables.py
\begin{table}[t]
  \centering
  \small
  \begin{tabular}{lrrrrr}
    \toprule
    Policy & capacity (s) & watermarked (s) & ordinary (s) & detection (s) & bases/s \\
    \midrule
    \policyCtok & 126.2 & 198.2 & 174.3 & 53.3 & 124.0 \\
    \policyGtok & 268.3 & 645.8 & 630.2 & 57.9 & 38.1 \\
    \policyGbp & 279.8 & 677.4 & 660.7 & 53.8 & 36.3 \\
    \bottomrule
  \end{tabular}
  \caption{Runtime envelope on one laptop, from single observations under uncontrolled desktop load and on mains power: a feasibility statement, not a benchmark. Detection uses neither the model nor the GPU. No peak-memory figure is reported, because the recorded memory counters cannot support a residency claim. The same generation runs roughly ten times slower on battery.}
  \label{tab:runtime}
\end{table}

